\chapter{Parallel \& Sequential Loops}
\label{chap:loops_forall}

\section{Parallel Loops (\texttt{for})}
The parallel \texttt{for} loop (lowered to \texttt{ForallNode} compound graphs in IF1 IR) is the primary construct for data-parallel execution in Sisal-2026. Every iteration of a parallel \texttt{for} loop is independent and can execute concurrently on separate CPU worker threads or GPU thread blocks.

\begin{note}[Keyword Clarification: \texttt{for} vs. \texttt{forall}]
In Sisal-2026 source code, the keyword for parallel loops is \textbf{\texttt{for}} (e.g., \texttt{for i in 1..N ... end for}). The term \emph{forall} refers to the underlying IF1 dataflow compound graph node (\texttt{ForallNode}) and literature terminology.
\end{note}

\subsection{Syntax \& Generator Ranges}
A parallel \texttt{for} loop consists of three sections:
\begin{enumerate}
    \item \textbf{Generator Range}: Defines the index domain (\texttt{i in 1, N} or element iteration \texttt{x in Array}).
    \item \textbf{Loop Body}: Computes expression bindings for each element in the domain.
    \item \textbf{Returns Clause}: Specifies reduction accumulators (\texttt{sum of}, \texttt{array of}, \texttt{stream of}, \texttt{product of}, \texttt{least of}, \texttt{greatest of}).
\end{enumerate}

\begin{lstlisting}[caption={Parallel \texttt{for} Loop Example}]
function VectorScale( A : array_dv[real]; factor : real returns array_dv[real] )
  for x in A
    y := x * factor
  returns
    array of y
  end for
end function
\end{lstlisting}

\subsection{Cross Products \& Dot Products}
Multiple generator ranges can be combined using \texttt{cross} (Cartesian product) or \texttt{dot} (lock-step parallel iteration):

\begin{lstlisting}[caption={Cross Product vs. Dot Product Generators}]
% Cross Product (N x M iterations)
Matrix := for i in 1, N cross j in 1, M
  val := i * 10 + j
returns
  array of val
end for;

% Dot Product (min(N, M) lock-step iterations)
Sums := for x in VectorA dot y in VectorB
  z := x + y
returns
  array of z
end for;
\end{lstlisting}

\section{Sequential Loops (\texttt{for initial})}
When loop iterations have loop-carried dependencies (where iteration $k$ requires results from iteration $k-1$), the \texttt{for initial} sequential loop is used.

\begin{lstlisting}[caption={Sequential Loop with Zero-Copy Tuple Swapping}]
function Fibonacci( n : integer returns integer )
  for initial
    i := 1;
    a := 0;
    b := 1
  while i < n repeat
    i := old i + 1;
    a := old b;
    b := old a + old b
  returns
    value of b
  end for
end function
\end{lstlisting}

\begin{note}[\texttt{old} Keyword]
Inside a \texttt{for initial} loop body, the \texttt{old} keyword references value bindings from the previous iteration. The compiler lowers state updates to zero-copy pointer swaps.
\end{note}

\section{Loop Semantics \& Comparative Analysis: Sisal vs. C vs. Fortran}
Sisal's \texttt{for initial} construct is a functional state-transition graph loop rather than a control-flow jump. This fundamental semantic difference distinguishes Sisal loops from traditional imperative loops in C and Fortran.

\subsection{Core Rules of \texttt{for initial} Iteration}
\begin{enumerate}
    \item \textbf{Rule 1: Unconditional Initial Seed Evaluation (No Undefined States)}:
    Sisal defines a single-assignment value history that begins with an explicit initial state (\texttt{body\_0}). The \texttt{initial} seed is unconditionally evaluated and added to any accumulation collection (e.g. \texttt{returns stream of I}) \emph{before} the loop condition check occurs. Even in a zero-trip loop where the condition is false on entry, the initial seed value is safely defined and gathered, eliminating any risk of reading garbage or uninitialized memory.
    
    \item \textbf{Rule 2: Out-of-Bounds Value Calculation \& Gathering}:
    In Sisal, the loop body computes the new state of loop variables from the previous state (\texttt{I := old I + step}). If the condition is satisfied, the body computes the new state and immediately gathers it. Consequently, the final value of \texttt{I} (which fails the subsequent condition check and terminates the loop) has already been computed and gathered into the result collection.
\end{enumerate}

\subsection{Comparative Analysis: Zero-Trip \& Normal Loops}
Table~\ref{tab:loop_comparison_zero_trip} and Table~\ref{tab:loop_comparison_normal} compare the execution behavior of Sisal, C/C++, and Fortran under zero-trip and normal loop scenarios.

\begin{table}[h!]
\centering
\small
\begin{tabularx}{\textwidth}{l >{\raggedright\arraybackslash}p{4.2cm} >{\raggedright\arraybackslash}X}
\toprule
\textbf{Language} & \textbf{Example Syntax} & \textbf{Execution Details \& Undefined Risk} \\
\midrule
\textbf{Sisal-2026} & \texttt{for initial I := 10 while I < 5 repeat I := old I + 1 returns stream of I end for} & Initial seed executes unconditionally and gathers \texttt{10}. Condition \texttt{10 < 5} is false $\rightarrow$ exits. \\
\addlinespace
& & \textbf{Stream}: \texttt{[10]} (Size 1) \quad \textbf{Risk}: None. \\
\midrule
\textbf{C / C++} & \texttt{int r; for(int i=10; i<5; i++) r=i;} & Condition \texttt{10 < 5} is checked first and fails. Body never executes. \\
\addlinespace
& & \textbf{Vector}: \texttt{[]} (Empty) \quad \textbf{Risk}: High (reading \texttt{r} yields garbage/UB). \\
\midrule
\textbf{Fortran} & \texttt{INTEGER :: I, R; DO I = 10, 5; R = I; END DO} & Trip count \texttt{MAX(0, 5-10+1)} evaluates to \texttt{0}. Body never executes. \\
\addlinespace
& & \textbf{Result}: Uninitialized \quad \textbf{Risk}: High. \\
\bottomrule
\end{tabularx}
\caption{Zero-Trip Loop Behavior Comparison (Initial Seed = 10, Bound = 5).}
\label{tab:loop_comparison_zero_trip}
\end{table}

\begin{table}[h!]
\centering
\small
\begin{tabularx}{\textwidth}{l >{\raggedright\arraybackslash}p{4.2cm} >{\raggedright\arraybackslash}X}
\toprule
\textbf{Language} & \textbf{Example Syntax} & \textbf{State Sequence \& Resulting Collection} \\
\midrule
\textbf{Sisal-2026} & \texttt{for initial I := 1 while I <= 5 repeat I := old I + 1 returns stream of I end for} & 1. Seed gathers \texttt{1}. \newline 2. Iterations 1--5 compute and gather \texttt{2, 3, 4, 5, 6}. \newline 3. \texttt{6 <= 5} fails $\rightarrow$ exits. \\
\addlinespace
& & \textbf{Stream}: \texttt{[1, 2, 3, 4, 5, 6]} (Size 6) \\
\midrule
\textbf{C / C++} & \texttt{for(int i=1; i<=5; i++) vec.push\_back(i);} & 1. Iterations 1--5 push \texttt{1, 2, 3, 4, 5}. \newline 2. Increment to \texttt{6}; \texttt{6 <= 5} fails $\rightarrow$ exits. \\
\addlinespace
& & \textbf{Vector}: \texttt{[1, 2, 3, 4, 5]} (Size 5) \\
\midrule
\textbf{Fortran} & \texttt{DO I = 1, 5; ... END DO} & Loop executes 5 times. Final \texttt{I} is incremented to \texttt{6} on exit. \\
\bottomrule
\end{tabularx}
\caption{Normal 5-Iteration Loop Behavior Comparison (Range = 1..5).}
\label{tab:loop_comparison_normal}
\end{table}

\subsection{Key Design Takeaways}
\begin{itemize}
    \item \textbf{Collection Sizes}: For $N$ successful condition evaluations, a traditional C/Fortran loop produces $N$ elements, whereas a Sisal \texttt{for initial} loop yields $N + 1$ elements (the initial seed plus $N$ body calculations).
    \item \textbf{Initialization Safety}: Sisal guarantees that loop-carried value streams always contain at least the initial seed value, statically eliminating uninitialized memory access.
\end{itemize}
